\chapter{Massive Particles and Spin}
\label{ch:massive-particles-spin}

\section{Massive One-Particle Representations}

This chapter works in \(D=4\). Let \(\Hilb_1\) be an irreducible massive
one-particle subspace of mass \(m>0\). The positive-energy mass shell is
\[
  \Sigma_m^+
  =
  \{p=(p^0,\vec p):p^0=\omega_{\vec p},\ p^2=-m^2\},
  \qquad
  \omega_{\vec p}=\sqrt{\vec p^{\,2}+m^2}.
\]
Generalized momentum states are normalized by
\[
  \langle \vec q,\tau|\vec p,\sigma\rangle
  =
  (2\pi)^3\,2\omega_{\vec p}\,
  \delta^3(\vec q-\vec p)\,\delta_{\tau\sigma}.
\]
The labels \(\sigma,\tau\) span a finite-dimensional internal vector space
which is fixed by Lorentz symmetry.

Choose the reference momentum
\[
  p_R=(m,\vec0).
\]
For every \(p\in\Sigma_m^+\), choose a standard Lorentz transformation
\[
  L(p)\in SO^+(1,3),
  \qquad
  L(p)p_R=p.
\]
The stabilizer of \(p_R\) is the rotation group \(SO(3)\). Its double cover
\[
  \operatorname{Spin}(3)\simeq SU(2)
\]
acts on quantum states because Hilbert-space symmetries act on rays and may
therefore be projective before lifting to the simply connected cover.
In state-space formulas below the Lorentz transformations are understood as
elements of the universal cover when they are inserted into spin
representations; the same letters denote their projections to
\(SO^+(1,3)\) when they act on four-vectors.

For a Lorentz transformation \(\Lambda\), define the Wigner rotation
\[
  W(\Lambda,p)
  =
  L(\Lambda p)^{-1}\Lambda L(p).
\]
It obeys
\[
  W(\Lambda,p)p_R=p_R,
\]
and is therefore an element of the massive little group. If the particle has
spin \(j\), with
\[
  j=0,\frac12,1,\frac32,\ldots,
\]
let \(D^{(j)}\) be the corresponding irreducible unitary representation of
\(SU(2)\) on a vector space \(V_j\) of dimension \(2j+1\). The one-particle
states transform as
\[
  U(\Lambda)|\vec p,\sigma\rangle
  =
  \sum_{\tau=-j}^{j}
  D^{(j)}_{\tau\sigma}\!\left(W(\Lambda,p)\right)
  |\Lambda\vec p,\tau\rangle ,
\]
where \(\sigma,\tau=-j,-j+1,\ldots,j\). The covariant normalization above
absorbs the Jacobian of the mass-shell measure, so no additional
\(\sqrt{(\Lambda p)^0/p^0}\) factor appears.

\begin{center}
\begin{tikzpicture}[scale=0.95,>=Stealth]
  \draw[blue!70!black,thick] (-3.0,0) ellipse (1.55 and 1.05);
  \node[blue!70!black] at (-3.0,1.35) {$SO(3)$};
  \fill (-3.75,0) circle (2pt) node[left] {\(1\)};
  \fill (-2.25,0) circle (2pt) node[right] {\(2\pi\)};
  \draw[purple!80!black,very thick,->]
    (-3.65,-0.14).. controls (-3.35,-0.72) and (-2.65,-0.72) .. (-2.35,-0.14);
  \node[purple!80!black,fill=white,inner sep=1.5pt] at (-3.0,-1.18)
    {closed loop};

  \draw[->,gray!70,very thick] (-0.85,0)--(0.85,0)
    node[midway,above] {lift};

  \draw[green!50!black,thick] (3.0,0) ellipse (1.75 and 1.12);
  \node[green!50!black] at (3.0,1.38) {$SU(2)$};
  \fill (2.1,0) circle (2pt) node[left] {\(1\)};
  \fill (3.9,0) circle (2pt) node[right] {\(-1\)};
  \draw[red!75!black,very thick,->]
    (2.25,-0.15).. controls (2.7,-0.72) and (3.3,-0.72) .. (3.75,-0.15);
  \node[red!75!black,fill=white,inner sep=1.5pt] at (3.0,-0.94) {\(2\pi\)};
  \draw[orange!80!black,very thick,->]
    (3.75,0.15).. controls (3.3,0.72) and (2.7,0.72) .. (2.25,0.15);
  \node[orange!80!black,fill=white,inner sep=1.5pt] at (3.0,0.94) {\(4\pi\)};
\end{tikzpicture}
\end{center}

A \(2\pi\) rotation acts on \(V_j\) by
\[
  (-1)^{2j}.
\]
Integer and half-integer spins are therefore distinct representations of the
same massive little-group construction.

\section{Spin Frames on the Mass Shell}

The choice of standard transformations \(L(p)\) is a choice of spin frame over
the mass shell. To make this statement precise, let
\[
  d\mu_m(p)=\frac{d^3p}{(2\pi)^3\,2\omega_{\vec p}}
\]
be the invariant measure on \(\Sigma_m^+\). After choosing \(L(p)\), the
one-particle Hilbert space is represented as
\[
  \Hilb_1\simeq L^2(\Sigma_m^+,d\mu_m;V_j).
\]
The generalized basis associated to this frame is the distributional map
\[
  F_L(p):V_j\longrightarrow \Hilb_1,
  \qquad
  F_L(p)e_\sigma=|\vec p,\sigma\rangle ,
\]
where \(\{e_\sigma\}_{\sigma=-j}^j\) is an orthonormal basis of \(V_j\).
The Lorentz action is
\[
  U(\Lambda)F_L(p)
  =
  F_L(\Lambda p)D^{(j)}(W(\Lambda,p)).
\]
This formula is a representation of the Lorentz group because the Wigner
rotation obeys the cocycle identity
\[
  W(\Lambda_2\Lambda_1,p)
  =
  W(\Lambda_2,\Lambda_1 p)W(\Lambda_1,p).
\]
Indeed,
\[
\begin{aligned}
  W(\Lambda_2\Lambda_1,p)
  &=
  L(\Lambda_2\Lambda_1p)^{-1}\Lambda_2\Lambda_1L(p)       \\
  &=
  L(\Lambda_2\Lambda_1p)^{-1}\Lambda_2L(\Lambda_1p)\,
  L(\Lambda_1p)^{-1}\Lambda_1L(p).
\end{aligned}
\]

If \(Q:\Sigma_m^+\to SU(2)\) is a smooth little-group-valued function, then
\[
  L'(p)=L(p)Q(p)
\]
is another choice of standard transformations. The corresponding Wigner
rotation is
\[
  W'(\Lambda,p)
  =
  Q(\Lambda p)^{-1}W(\Lambda,p)Q(p).
\]
The associated generalized basis is
\[
  F_{L'}(p)=F_L(p)D^{(j)}(Q(p)).
\]
Thus a wavefunction written as a column \(\psi_L(p)\in V_j\) in the \(L\)
frame is written in the \(L'\) frame as
\[
  \psi_{L'}(p)=D^{(j)}(Q(p))^{-1}\psi_L(p).
\]
The inner product
\[
  \langle \psi,\chi\rangle
  =
  \int_{\Sigma_m^+} d\mu_m(p)\,
  \psi_L(p)^\dagger\chi_L(p)
\]
is independent of the spin frame because \(D^{(j)}(Q(p))\) is unitary. In this
language, a scattering amplitude with external massive spin states is a
multilinear pairing of spin-frame components and dual spin-frame components
whose value is invariant under the simultaneous transformations above.

\section{Covariant Fields as Intertwiners}

Let \(\widehat\Phi_A(x)\) be a local field whose index \(A\) transforms in a
finite-dimensional representation \(R\) of the double cover
\(\operatorname{Spin}^+(1,3)\simeq SL(2,\mathbb C)\):
\[
  U(\Lambda)\widehat\Phi_A(x)U(\Lambda)^{-1}
  =
  R_A{}^B(\Lambda)\widehat\Phi_B(\Lambda x).
\]
The vacuum-to-particle matrix element is
\[
  \langle\Omega|\widehat\Phi_A(0)|\vec p,\sigma\rangle
  =
  u_A{}^\sigma(p).
\]
Translation covariance gives the spacetime dependence
\[
  \langle\Omega|\widehat\Phi_A(x)|\vec p,\sigma\rangle
  =
  \ee^{\ii p\cdot x}u_A{}^\sigma(p).
\]
Lorentz covariance and the Wigner transformation law imply
\[
  \sum_{\tau=-j}^j
  u_A{}^\tau(\Lambda p)
  D^{(j)}_{\tau\sigma}(W(\Lambda,p))
  =
  R_A{}^B(\Lambda^{-1})u_B{}^\sigma(p).
\]
This is the equivariance condition that connects the covariant field index
\(A\) to the spin-frame index \(\sigma\).

At the reference momentum \(p_R\), the matrix
\[
  u_A{}^\sigma(p_R):V_j\longrightarrow R|_{SU(2)}
\]
is an intertwiner of \(SU(2)\) representations:
\[
  R(W)_A{}^B u_B{}^\sigma(p_R)
  =
  \sum_{\tau=-j}^j
  u_A{}^\tau(p_R)D^{(j)}_{\tau\sigma}(W),
  \qquad W\in SU(2).
\]
Thus a covariant field can have nonzero overlap with a spin-\(j\) particle
exactly through the spin-\(j\) part contained in the restriction of \(R\) to
the massive little group. The boosted wavefunctions are
\[
  u_A{}^\sigma(p)
  =
  R_A{}^B(L(p)^{-1})\,u_B{}^\sigma(p_R),
\]
with the convention for \(L(p)\) fixed above. Different choices of standard
boost are spin-frame changes of the type described in the previous section,
and the field intertwiners are used with the corresponding transformed
spin-frame components.

\section{Spin One Half}

For spin \(j=\frac12\), the smallest Lorentz-covariant field space is a Dirac
spinor. Let the gamma matrices obey
\[
  \{\gamma^\mu,\gamma^\nu\}=2\eta^{\mu\nu},
  \qquad
  \eta=\operatorname{diag}(-,+,+,+).
\]
Define
\[
  S^{\mu\nu}
  =
  -\frac{\ii}{4}[\gamma^\mu,\gamma^\nu],
  \qquad
  R(\Lambda)
  =
  \exp\!\left(-\frac{\ii}{2}\omega_{\mu\nu}S^{\mu\nu}\right)
\]
for
\[
  \Lambda=\exp(\omega),
  \qquad
  \omega_{\mu\nu}=-\omega_{\nu\mu}.
\]
The generators satisfy
\[
  [S^{\mu\nu},\gamma^\rho]
  =
  -\ii\gamma^\mu\eta^{\nu\rho}
  +\ii\gamma^\nu\eta^{\mu\rho},
\]
and therefore
\[
  R(\Lambda)\gamma^\mu R(\Lambda)^{-1}
  =
  \Lambda^\mu{}_\nu\gamma^\nu .
\]

Write
\[
  \not p=p_\mu\gamma^\mu .
\]
For \(p_R=(m,\vec0)\),
\[
  \not p_R=-m\gamma^0,
  \qquad
  \not p_R^{\,2}=-m^2 .
\]
Choose two rest spinors \(u_R^\sigma\) and two rest spinors \(v_R^\sigma\),
\(\sigma=1,2\), by
\[
  \not p_R u_R^\sigma=\ii m\,u_R^\sigma,
  \qquad
  \not p_R v_R^\sigma=-\ii m\,v_R^\sigma .
\]
Their boosts are
\[
  u^\sigma(p)=R(L(p)^{-1})u_R^\sigma,
  \qquad
  v^\sigma(p)=R(L(p)^{-1})v_R^\sigma .
\]
Using covariance of \(\gamma^\mu\), these spinors obey
\[
  (\not p-\ii m)u^\sigma(p)=0,
  \qquad
  (\not p+\ii m)v^\sigma(p)=0.
\]
These are the momentum-space Dirac equations associated with the two
eigenspaces of \(\not p\) on the mass shell.

Let
\[
  \beta=\ii\gamma^0,
  \qquad
  \bar u=u^\dagger\beta,
  \qquad
  \bar v=v^\dagger\beta .
\]
The projectors onto the two eigenspaces of \(\not p\) are
\[
  \Pi_+(p)=\frac{m-\ii\not p}{2m},
  \qquad
  \Pi_-(p)=\frac{m+\ii\not p}{2m}.
\]
With an orthonormal basis of rest spinors, the spin sums are proportional to
these projectors:
\[
  \sum_{\sigma=1}^2 u^\sigma(p)\bar u^\sigma(p)
  \propto
  m-\ii\not p,
  \qquad
  \sum_{\sigma=1}^2 v^\sigma(p)\bar v^\sigma(p)
  \propto
  m+\ii\not p,
\]
where the proportionality constant is fixed by the normalization chosen for
one-particle spinor wavefunctions. These projector identities are the algebraic
source of the external spinor factors in spinorial LSZ reduction.

\section{Particle Spin and Field Indices}

The spin \(j\) labels an irreducible unitary representation of the massive
little group on the one-particle Hilbert space. A field index \(A\) labels a
finite-dimensional Lorentz representation used by a local operator. The two
are related by the intertwiner
\[
  u_A{}^\sigma(p)
  =
  \langle\Omega|\widehat\Phi_A(0)|\vec p,\sigma\rangle .
\]
This distinction is part of the basic architecture of relativistic quantum
theory. The same spin-\(j\) particle may be created by different local
operators, and a single covariant field representation may contain several
little-group spins when restricted to \(SU(2)\). LSZ for spinning particles
uses the pole of the relevant two-point function together with the
intertwiner \(u_A{}^\sigma(p)\) to convert field indices into external
particle spin labels.

The next chapter constructs spinor field operators and their functional
integral calculus. The representation-theoretic data established here will
then become the spinor completeness relations, propagator numerators, and
external-line factors used in perturbative calculations.
